\section{The local estimate and the coefficient algebra}\label{sec:local}

The analytic estimate concerns four inputs on the real line. Its
constant may depend on the slopes, but must be independent of the number
of scales and of the number of changes of the inputs. This uniformity
will give the oscillation estimate in Section~\ref{sec:transfer}.

\subsection{Normalization and local forms}
Set
\[
 g=\gcd(\abs{\lambda_1},\abs{\lambda_2},\abs{\lambda_3}),
 \qquad S=T^g,\qquad b_0=0,\qquad b_j=\lambda_j/g\quad(1\le j\le3).
\]
Then $S$ is invertible, probability preserving, and mixing, and
$T^{\lambda_j n}=S^{b_jn}$. Henceforth the four distinct integers
$b_0,b_1,b_2,b_3$ are fixed, and
$\gcd(\abs{b_1},\abs{b_2},\abs{b_3})=1$. We make no assertion of
uniformity as these integers vary.

Dyadic intervals are half-open. The depth of a descendant $I$ of a
dyadic interval $J$ is $s(I)=\log_2(\abs{J}/\abs{I})$. A finite rooted
tree contains every ancestor between any of its intervals and its root.
The dyadic lattice may be translated.

\begin{definition}[Local kernels]\label{def:kernels}
Fix $\delta_0>0$ and $C_0\ge1$. For an interval $I$ of length $r$, a kernel
$w_I\in C_c^\infty(\R^2;\C)$ is admissible if
\begin{enumerate}
\item $\dist(x+b_jt,\partial I)\ge\delta_0r$ and $x+b_jt\in I$ for every
$(x,t)\in\supp w_I$ and $0\le j\le3$;
\item for every $v\in\R$ and $0\le j\le3$,
\begin{equation}\label{eq:marginals}
 \int_\R w_I(v-b_jt,t)\,dt=0;
\end{equation}
\item for nonnegative integers $a,d$ and $n=a+d$,
\begin{equation}\label{eq:derivatives}
 \norm{\partial_x^a\partial_t^d w_I}_\infty
 \le r^{-n}\bigl(C_0(n+2)\bigr)^{C_0(n+2)}.
\end{equation}
\end{enumerate}
For scalar measurable inputs define
\begin{equation}\label{eq:local-form}
 H_I(z_0,z_1,z_2,z_3)
 =r^{-2}\int_{\R^2} w_I(x,t)\prod_{j=0}^3z_j(x+b_jt)\,dx\,dt,
\end{equation}
whenever the integral is absolutely convergent.
\end{definition}

For a nonempty integer block $B=\{u,\ldots,v\}$ and functions $z_s$ on
a common domain, put
\[
 \mathcal V_B(z)=\max_{u\le s\le v}\abs{z_s}
       +\sum_{s=u}^{v-1}\abs{z_{s+1}-z_s}.
\]
This is a pointwise bound on both size and variation within the block;
no variation across block boundaries is included.

\begin{theorem}[Local estimate with variation in depth]\label{thm:local}
Fix the primitive distinct integers $b_0=0,b_1,b_2,b_3$ and the
parameters $\delta_0,C_0$. Let $\mathcal T$ be a finite dyadic tree with root
$J$ and largest depth $d$, and let $\mathcal I\subseteq\mathcal T$.
Assign an admissible kernel to every $I\in\mathcal I$.
For each $0\le j\le3$, let $z_{j,s}:J\to\C$, $0\le s\le d$, be
measurable, and partition $\{0,\ldots,d\}$ into $W_j\ge1$
deterministic nonempty consecutive blocks. Suppose that
\[
 \mathcal V_B(z_j)\le W_j^{-1/2}\quad\text{almost everywhere on }J
\]
for each block in the $j$th partition. Extend every input by zero
outside $J$. Then
\begin{equation}\label{eq:local-estimate}
 \sum_{I\in\mathcal I}\abs{I}\,
 \abs{H_I(z_{0,s(I)},z_{1,s(I)},z_{2,s(I)},z_{3,s(I)})}
 \le C_{b,\delta_0,C_0}\abs{J}.
\end{equation}
The constant is independent of the lattice, tree, depth range,
subcollection, kernels, and the four partitions and their sizes.
\end{theorem}

The schedules are prescribed at every depth, even depths with no
summation intervals. This is exactly the hypothesis of
\cite[Theorem~3.2]{V}, except for the coefficients in the local forms.
Sections~\ref{sec:analytic}--\ref{sec:completion} prove the theorem by
extending that proof and its analytic inputs from~\cite{H}. Before
turning to those technical arguments, Section~\ref{sec:transfer} proves
the implication from Theorem~\ref{thm:local} to the ergodic theorem.

\subsection{Integer reconstruction and quadratic cancellation}
Two features of the coefficients recur in both the analytic proof and
transference. Pairs of coordinates determine $(x,t)$ over $\R$, but
need not determine it over $\Z$. All four coordinates together do
determine the integral step. Quadratic cancellation requires a separate
identity, furnished by Vandermonde interpolation.

\begin{lemma}[Coefficient identities]\label{lem:algebra}
There exist integers $e_0,e_1,e_2,e_3$ and nonzero integers
$c_0,c_1,c_2,c_3$ such that
\begin{equation}\label{eq:bezout}
 \sum_{j=0}^3e_j=0,\qquad \sum_{j=0}^3e_jb_j=1,
\end{equation}
and
\begin{equation}\label{eq:vandermonde}
 \sum_{j=0}^3c_jb_j^a=0\qquad(a=0,1,2).
\end{equation}
The nullspace over $\R$ of the two rows $(1)_j,(b_j)_j$ is spanned
by $(c_j)_j$ and $(b_jc_j)_j$.
Consequently, for $y_j=x+b_jt$,
\[
 x=y_0,\qquad t=\sum_je_jy_j,\qquad \sum_jc_jy_j^2=0.
\]
For any distinct $i,j$, the pair map
$(x,t)\mapsto(x+b_it,x+b_jt)$ is injective on $\Z^2$ and pushes
normalized Haar measure on $(\R/L\Z)^2$ to normalized Haar measure,
for every $L>0$.
\end{lemma}

\begin{proof}
Bezout's identity gives integers $e_1,e_2,e_3$ with
$\sum_{j=1}^3e_jb_j=1$; set $e_0=-\sum_{j=1}^3e_j$.
Let $P(z)=\prod_{j=0}^3(z-b_j)$. Lagrange interpolation of $z^a$,
$0\le a\le2$, followed by comparison of the coefficient of $z^3$,
gives
$\sum_jb_j^a/P'(b_j)=0$. Multiply by a common positive integer
multiple of the denominators to obtain~\eqref{eq:vandermonde} with
every $c_j\ne0$. The two displayed null vectors are independent
because the $b_j$ are distinct, and the nullspace has dimension two.
The identities for $y_j$ follow by expansion. Finally the pair matrix
has determinant $b_j-b_i\ne0$, so it is injective on the lattice and
surjective on the real torus. A surjective continuous homomorphism of
compact groups pushes normalized Haar measure to normalized Haar
measure.
\end{proof}

\begin{remark}\label{rem:nonunimodular}
Primitivity does not mean that some pair is unimodular. For example,
the slopes $(b_0,b_1,b_2,b_3)=(0,6,10,15)$ are primitive, but no pair
difference is $1$ in absolute value. This is why the integer inverse
in~\eqref{eq:bezout}, rather than division by a pair determinant,
is used in the lattice and torus arguments below.
\end{remark}

\subsection{The fixed-input constant}
The proof of Theorem~\ref{thm:local} first establishes a fixed-input
estimate. For $2<q<3$, let $M_q f=(M(\abs f^q))^{1/q}$, where $M$
is the uncentered Hardy--Littlewood maximal operator. Let $C_N(q)$
be the least constant for which
\begin{equation}\label{eq:CN}
 \sum_{I\in\mathcal A}\abs I\,\abs{H_I(f_0,f_1,f_2,f_3)}
 \le C_N(q)\int_\R\prod_{j=0}^3M_qf_j(x)\,dx
\end{equation}
for all finite collections $\mathcal A$ of intervals in one translated
dyadic lattice, whose lengths belong to at most $N$ consecutive scales,
admissible kernels, and compactly supported
scalar $L^q$ inputs. This is the definition in
\cite[Section~2]{H}, with our coefficients.
Pair-coordinate estimates imply $C_N(q)\lesssim_bN$, so this constant
is finite before any absorption argument. The analytic proof will show
$\sup_NC_N(q)<\infty$ for a suitable fixed $q>2$, and then use that
bound in the scheduled-input argument. This order avoids using
Theorem~\ref{thm:local} to prove its own analytic input.
